\clearpage
\SourceLeaf{242}\pdfbookmark[1]{Section Three.}{liber-secundus-sectio-tertia}
\begin{center}
{\large SECTION THREE.}\par\medskip
\emph{Determination of an orbit satisfying any number of observations as nearly as possible.}\par\bigskip
\hypertarget{native-242-article172}{}\pdfbookmark[2]{172.}{article-172}172.
\end{center}
If the astronomical observations and the other numbers on which the computation of orbits rests had absolute precision, the elements, whether based on three observations or on four, would likewise be obtained immediately with absolute exactness (insofar, of course, as the motion is assumed to follow \textsc{Kepler}'s laws exactly), and so further observations could only confirm them, never correct them. But since all our measurements and observations are only approximations to the truth, and the same must hold for all computations based on them, the highest aim of all computations concerning concrete phenomena must be to approach the truth as closely as possible. This can be achieved only by suitably combining \emph{more} observations than are strictly required to determine the unknown quantities. We can undertake this task only once an approximate orbit has been found. It must then be corrected so that it satisfies all the observations \emph{as exactly as possible}. Although this expression may seem somewhat vague, we will give below the principles by which the problem can be solved rigorously and methodically.

It can be worth seeking the highest precision only when we are, so to speak, putting the finishing touches on the determination of the orbit. On the other hand, as long as there is hope that new observations will soon give \SourceLeaf{243}occasion for further corrections, it is advisable to relax the demand for the utmost precision to a greater or lesser extent according to circumstances, if this considerably shortens the computations. We shall try to provide for both cases.
